\documentclass[11pt]{article}
\usepackage[margin=1in]{geometry}
\usepackage{amsmath,amssymb,amsfonts,bm}
\usepackage{booktabs}
\usepackage{array}
\usepackage{enumitem}
\usepackage{xcolor}
\usepackage[colorlinks=true,linkcolor=blue!50!black,urlcolor=blue!50!black]{hyperref}

\newcommand{\code}[1]{\texttt{\detokenize{#1}}}
\DeclareMathOperator{\softplus}{softplus}
\DeclareMathOperator{\softmax}{softmax}
\DeclareMathOperator{\BCE}{BCE}

\title{Training the continuous-saturation emulator in nGibbs}
\author{Ben Thyer}
\date{\today}

\begin{document}
\sloppy
\maketitle

\begin{abstract}
This note describes how the continuous-saturation emulator (\code{ContinuousModel}) is
trained, with an emphasis on how each term of the loss is evaluated on a mini-batch. The
model is defined in \code{NN_continuous.py} and trained by \code{train_Upper_MELTS} in
\code{trainer.py}. Table~\ref{tab:key} is the key to all symbols.
\end{abstract}

%=====================================================================================
\section{Variables}
%=====================================================================================

Dimensions use the letter codes of the ML indexer (\code{src/ngibbs/config/README_MLIndexer.md}):
P phases, VP compositionally variable phases, C components (endmembers of all phases),
VC components of the variable phases, E elements and O oxides. Two further sizes are
used here: F, the number of physical input features (e.g.\ pressure and temperature),
and B, the number of rows in a mini-batch. Every quantity in Table~\ref{tab:key} is per
row; a batch stacks B rows. For emulator outputs, the hatted symbol is the prediction and
the plain symbol is the ground truth from the thermodynamic model.

\begin{table}[!tbp]
\centering
\small
\renewcommand{\arraystretch}{1.25}
\begin{tabular}{@{}c c c >{\raggedright\arraybackslash}p{10.2cm}@{}}
\toprule
Prediction & Truth & Dim. & Meaning \\
\midrule
\multicolumn{4}{@{}l}{\emph{Inputs}}\\
 & $u$ & F & Physical state (pressure, temperature, \dots), min--max normalised to $[0,1]$ \\
\midrule
\multicolumn{4}{@{}l}{\emph{Emulator outputs}}\\
$\hat b$ & $b$ & E & Bulk composition in element moles, normalised to sum to 1. $b$ is also a network input; $\hat b$ is the bulk reconstructed from the predicted phases \\
$\hat n$ & $n$ & P & Phase abundance (moles); zero for an absent phase \\
$\hat x$ & $x$ & VC & Endmember mole fractions of each variable phase (sum to 1 within a phase) \\
$\hat m$ & $m$ & C & Component (endmember) moles, $m = (\text{phase abundance}) \times (\text{fraction})$ \\
$\hat s$ & $s$ & P & Phase presence, 1 if present and 0 if absent; $\hat s = [\hat n > 0]$ \\
\midrule
\multicolumn{4}{@{}l}{\emph{Quantities used only in training}}\\
$g$ & -- & P & Phase affinity: the raw output of each phase's mole head; $\hat n = \max(g,0)$ \\
$\tilde n$ & -- & P & Smoothed abundance regressed by the mole loss, Eq.~\eqref{eq:smooth} \\
-- & $\bar s$ & VC & Presence $s$ copied to each endmember of its phase (chemistry mask) \\
\midrule
\multicolumn{4}{@{}l}{\emph{Training constants}}\\
\multicolumn{2}{c}{$T_0$} & P & Reference abundance scale per phase (5th percentile of present abundances) \\
\multicolumn{2}{c}{$T$} & P & Annealed boundary temperature, $T = a\,T_0$ \\
\multicolumn{2}{c}{$a$} & 1 & Annealing factor, decreased linearly over an episode \\
\multicolumn{2}{c}{$w$} & P & Phase weights (recipe \code{phase_weights}; default 1) \\
\multicolumn{2}{c}{$\omega$} & VC & Endmember weights (each endmember takes its phase's weight by default) \\
\multicolumn{2}{c}{$\lambda_s,\lambda_x,\lambda_n,\lambda_b$} & 1 & Weights of the four loss terms (recipe \code{loss_config}) \\
\midrule
\multicolumn{4}{@{}l}{\emph{Fixed stoichiometric matrices (from the ML indexer)}}\\
\multicolumn{2}{c}{$A$} & C$\times$E & Element moles per mole of each component (\code{compToEl}) \\
\multicolumn{2}{c}{$M$} & P$\times$C & Which components belong to which phase (\code{phaseToCompMap}) \\
\multicolumn{2}{c}{$V$} & VC$\times$C & Places the VC variable-phase fractions into the full component list (\code{variedToAllComp}) \\
\multicolumn{2}{c}{$f$} & C & Fixed fractions of stoichiometric phases (\code{fixed_phaseToCompMap}) \\
\bottomrule
\end{tabular}
\caption{Key to the symbols. ``--'' marks a quantity with no prediction or no ground-truth
counterpart.}
\label{tab:key}
\end{table}

%=====================================================================================
\section{The forward model}
%=====================================================================================

The network sees $[u, b]$. An encoder produces a feature vector $h$, which is concatenated
with the raw input. From it:
\begin{enumerate}[nosep]
  \item one softmax head per variable phase gives that phase's endmember fractions $\hat x$.
        Endmembers containing an element that is absent from $b$ are given zero probability;
  \item one small network per phase, which also sees $\hat x$, gives the affinity $g$;
  \item the abundance is the rectified affinity,
  \begin{equation}
    \hat n = \max(g, 0).
  \end{equation}
\end{enumerate}
Because $\max(g,0)$ is continuous at $g=0$, a phase leaving the assemblage is a kink in
$\hat n$, not a jump. (In the code $g$ first passes through a leaky ReLU and is then
clamped at zero; the result is identical.)

Component moles and the reconstructed bulk follow from stoichiometry:
\begin{align}
  \hat m &= (\hat x V + f) \odot (\hat n M), \label{eq:m}\\
  \hat b &= \frac{\hat m A}{\sum_e (\hat m A)_e}. \label{eq:bhat}
\end{align}
Here $\odot$ is the element-wise product, so $\hat m$ is each component's phase abundance
times its fraction in that phase. Mass balance (adjusting $\hat m$ until $\hat b = b$
exactly) is applied only at inference, by \code{NN_MELTS.apply_mass_balance}, and is not
part of the trained network.

%=====================================================================================
\section{The training objective}
%=====================================================================================

For each mini-batch the trainer minimises
\begin{equation}
\boxed{\;
  L = 10^4\big(\lambda_s L_s + \lambda_x L_x + \lambda_n L_n + \lambda_b L_b\big)
\;}
\label{eq:total}
\end{equation}
covering phase presence ($L_s$), chemistry ($L_x$), phase abundance ($L_n$) and bulk
($L_b$). The factor $10^4$ keeps the loss large compared with the optimiser's $\epsilon$.
The bulk term is added only when $\lambda_b \neq 0$. Typical recipe values are
$\lambda_x = 500$, $\lambda_n = 50$, $\lambda_s$ between $0.002$ and $0.02$, and
$\lambda_b = 0$ while tuning, rising to $50$ for the final joint training.
The same function (\code{_upper_loss}) computes the training loss and the validation
loss, so a model is always scored on the objective it was trained on.

In every term below, sums ($\sum$) run over all B rows of the batch as well as over the
index shown, and each term is a weighted mean: the weighted sum of errors divided by the
sum of the weights.

\subsection{The error function $\rho$}
Chemistry, abundance and bulk use the symmetric relative squared error
(\code{symmetric_rel_l2}),
\begin{equation}
  \rho(\hat y, y) = \frac{(\hat y - y)^2}{|\hat y| + |y|},
  \label{eq:rho}
\end{equation}
with the denominator floored at $10^{-6}$. The type is chosen separately for each term by
\code{loss_config.<term>.type} (\code{symmetric_rel_l1} and \code{huber} are also
available). Two limits explain its behaviour:
\begin{align}
  y = 0:\qquad & \rho(\hat y, 0) = |\hat y|, \label{eq:rho0}\\
  \hat y \approx y:\qquad & \rho(\hat y, y) \approx \frac{(\hat y - y)^2}{2y}. \label{eq:rhorel}
\end{align}
For a target of zero it acts like an absolute ($L_1$) error, a constant pull toward zero.
Near a non-zero target it is a squared error divided by the target, so a small abundance
counts as much as a large one in relative terms.

\subsection{Chemistry loss $L_x$}
Only phases that are present in the ground truth have a composition, so the predicted
fractions are masked by $\bar s$ before comparison and the mean runs over present
endmembers only:
\begin{equation}
  L_x = \frac{\sum_k \bar s_k\, \omega_k\, \rho(\bar s_k \hat x_k,\ x_k)}{\sum_k \bar s_k\, \omega_k},
  \qquad k = 1,\dots,\text{VC}.
\end{equation}

\subsection{Abundance loss $L_n$}
Every phase is supervised, including absent ones (target $n=0$): learning where an
abundance reaches zero is the point of the continuous model. The loss is not computed on
$\hat n$ itself but on a smoothed version of it controlled by the boundary temperature
$T$ (Section~\ref{sec:anneal}):
\begin{equation}
  \tilde n = T \softplus(g/T) = T \ln\!\big(1 + e^{g/T}\big),
  \label{eq:smooth}
\end{equation}
and
\begin{equation}
  L_n = \frac{\sum_j w_j\, \rho(\tilde n_j, n_j)}{\sum_j w_j},
  \qquad j = 1,\dots,\text{P}.
\end{equation}
$\tilde n$ is a smooth stand-in for $\max(g,0)$ with two useful properties:
\begin{equation}
  \frac{d\tilde n}{dg} = \operatorname{sigmoid}(g/T),
  \qquad
  0 < \tilde n - \max(g,0) \le T \ln 2 .
  \label{eq:smoothprops}
\end{equation}
The error relative to $\hat n$ is largest ($T\ln 2$) at $g = 0$, negligible once
$|g| > 3T$, and shrinks as $T$ is annealed. For an absent phase, Eq.~\eqref{eq:rho0} gives
$\rho(\tilde n, 0) = \tilde n$, so the gradient on its affinity is
$\operatorname{sigmoid}(g/T)$. This pushes $g$ negative and fades to zero once the phase
is confidently absent ($g \ll -T$), instead of holding it at the boundary $g=0$.

\subsection{Presence loss $L_s$}
There is no separate presence head. Instead the affinity is also used as the logit of a
binary cross-entropy against $s$:
\begin{equation}
  L_s = \frac{\sum_j c_j\, \BCE\!\big(5g_j/T_{0,j},\ s_j\big)}{\sum_j c_j},
  \qquad
  c = \big(s\,w + (1-s)\big)\,\big|\tanh(g/T)\big|,
  \label{eq:sat}
\end{equation}
with $\BCE(z,s) = -s\ln\operatorname{sigmoid}(z) - (1-s)\ln\big(1-\operatorname{sigmoid}(z)\big)$.
Three choices are built into Eq.~\eqref{eq:sat}:
\begin{enumerate}[nosep]
  \item Phase weights apply only to the present examples ($s=1$). Those are the examples a
        rare phase is short of; absent examples keep weight 1.
  \item The logit is scaled by the fixed $T_0$, not the annealed $T$, so the classifier's
        confidence means the same thing throughout training: it reaches
        $\operatorname{sigmoid}(5) \approx 0.99$ at $g = T_0$.
  \item The weight $|\tanh(g/T)|$ uses the annealed $T$. It is zero at the boundary and
        close to 1 a few $T$ away. It hands off with the abundance loss, whose gradient
        $\operatorname{sigmoid}(g/T)$ is largest near the boundary and fades far from it:
        $L_n$ handles calibration close to $g=0$, and $L_s$ corrects samples that are
        confidently wrong far from it.
\end{enumerate}
The weight $c$ depends on $g$ and is not detached, so it contributes to the gradient.

\subsection{Bulk loss $L_b$}
The reconstructed bulk (Eq.~\ref{eq:bhat}) is compared with the input bulk over the
elements present in it:
\begin{equation}
  L_b = \frac{\sum_{e:\, b_e > 0} \rho(\hat b_e, b_e)}{\sum_{e:\, b_e > 0} 1},
  \qquad e = 1,\dots,\text{E}.
\end{equation}
This measures how far the raw predictions are from satisfying mass balance, which reduces
the correction that has to be applied at inference. It is meaningless until the mole heads
have been trained, so it is switched on only for the final joint training. $\hat b$ is
built from the unsmoothed $\hat n$.

%=====================================================================================
\section{Boundary temperature}
\label{sec:anneal}
%=====================================================================================

\paragraph{The reference scale $T_0$.}
$T_0$ sets, for each phase, how small an abundance is ``small''. It is the 5th percentile
of that phase's abundance over the rows in which the phase is present:
\begin{equation}
  T_{0,j} = 5\text{th percentile of } \{\, n_j : n_j > 0 \,\}.
\end{equation}
An abundance below $T_{0,j}$ is therefore smaller than in 95\% of the rows where phase $j$
is present. Because the training and validation sets are generated the same way, their
$T_0$ values are essentially identical, so $T_0$ is computed on the smaller validation set,
where it is cheaper. Across a typical dataset $T_0$ spans about three orders of magnitude
(roughly $10^{-5}$ for trace phases to $0.06$ for dominant ones), which is why one
temperature for all phases would not work. A phase that never occurs gets $T_0 = 0$; for
the division in Eqs.~\eqref{eq:smooth}--\eqref{eq:sat} it is replaced by $10^{-6}$ times
the smallest non-zero $T_0$, which makes its smoothing negligible.

\paragraph{Annealing.}
During an episode of $N$ epochs, the temperature is $T = a\,T_0$, with $a$ decreasing
linearly from $a_{\text{start}}$ to $a_{\text{end}}$:
\begin{equation}
  a = a_{\text{start}} + (a_{\text{end}} - a_{\text{start}})\,\frac{\text{epoch}}{N-1},
  \qquad \text{epoch} = 0, \dots, N-1.
\end{equation}
Early in training a large $T$ gives every phase a wide, smooth transition region and
well-behaved gradients. As $T$ shrinks, $\tilde n$ approaches $\hat n = \max(g,0)$ and the
model is trained on its actual output. The schedule restarts in every episode that enables
it (\code{boundary_temperature} in the recipe). Recipes typically chain two episodes,
$a: 1 \to 0.05$ and then $a: 0.05 \to 10^{-4}$. All arithmetic involving $g/T$, and the mole
heads that produce $g$, run in 32-bit floats even under mixed precision, because late in
the schedule $g/T$ would overflow or lose precision in 16-bit formats.

%=====================================================================================
\section{The training loop}
%=====================================================================================

\paragraph{Episodes.}
A recipe is a sequence of episodes, each with its own number of epochs, learning-rate
schedule, frozen modules and loss-weight overrides. The continuous recipes run:
\begin{enumerate}[nosep]
  \item tune the encoder and chemistry heads, with the mole heads frozen and
        $\lambda_n = \lambda_s = 0$;
  \item tune the mole heads, with the encoder and chemistry heads frozen and
        $\lambda_x = 0$;
  \item train everything jointly, with annealing on and $\lambda_b > 0$.
\end{enumerate}

\paragraph{Optimiser and input noise.}
AdamW (with AMSGrad and $\epsilon = 10^{-4}$), with weight decay set per module: encoder
\code{lowWD}, chemistry heads \code{highWD}, mole heads none. Training inputs can be
perturbed with multiplicative Gaussian noise, $[u,b] \leftarrow [u,b]\,(1 + \sigma\xi)$
with $\xi \sim \mathcal N(0,1)$; the perturbed $b$ is also the bulk target in $L_b$.

\paragraph{Validation, checkpointing and early stopping.}
Before the first epoch, the incoming model is scored on the validation set; that score is
the one to beat. After each epoch the validation loss (Eq.~\ref{eq:total}, at that epoch's
$T$) is computed, and the model is saved whenever it improves on the best so far. Training
stops after \code{early_stopping_patience} epochs without improvement, and the best saved
model is restored.

\paragraph{Adaptive regularisation.}
After each epoch, if the validation loss exceeds 1.02 times the training loss, the
dropout rate (outside the mole heads) is raised one step up to its ceiling; once at the
ceiling, the input noise $\sigma$ is raised instead. If the validation loss is below the
training loss, dropout (or, once dropout is zero, $\sigma$) is lowered one step.

\paragraph{Remark on checkpoint selection under annealing.}
The validation loss is evaluated at a different $T$ each epoch. Both the smoothing error
($\le T\ln 2$) and $\tilde n$ for absent phases shrink with $T$, so the validation loss
tends to fall as the schedule progresses regardless of model quality. Checkpoint
selection and early stopping therefore favour later epochs, and the starting score
(computed at the first epoch's $T$) is not directly comparable with later ones.

\end{document}
